%% start of file `template.tex'.
%% Copyright 2006-2013 Xavier Danaux (xdanaux@gmail.com).
%
% This work may be distributed and/or modified under the
% conditions of the LaTeX Project Public License version 1.3c,
% available at http://www.latex-project.org/lppl/.

% possible options include font size ('10pt', '11pt' and '12pt'), paper size ('a4paper', 'letterpaper', 'a5paper', 'legalpaper', 'executivepaper' and 'landscape') and font family ('sans' and 'roman')
\documentclass[10pt,a4paper,sans]{maxbothe-cv}

% Remove dot at the end of a eventry
\usepackage{xpatch}
\xpatchcmd{\cventry}{.\strut}{\strut}{}{}

\newif\ifconfidential
\confidentialtrue  % Uncomment to hide confidential information
%\confidentialfalse  % Uncomment to show confidential information


% personal data
\name{Max}{Bothe}
%\title{}                       % optional, remove / comment the line if not wanted
\email{me@maxbothe.de}                         % optional, remove / comment the line if not wanted
\homepage{maxbothe.de}                            % optional, remove / comment the line if not wanted
\social[linkedin]{maxbothe}                         % optional, remove / comment the line if not wanted
\social[github]{mathebox}                          % optional, remove / comment the line if not wanted
\social[researchgate][https://www.researchgate.net/profile/Max-Bothe]{ Max Bothe}
%\extrainfo{additional information}                % optional, remove / comment the line if not wanted
\photo[90pt][0pt]{picture}                       % optional, remove / comment the line if not wanted; '64pt' is the height the picture must be resized to, 0.4pt is the thickness of the frame around it (put it to 0pt for no frame) and 'picture' is the name of the picture file
%\quote{Some quote}                                 % optional, remove / comment the line if not wanted

\ifconfidential
\else
	\input{confidential}
\fi


%----------------------------------------------------------------------------------
%            content
%----------------------------------------------------------------------------------
\begin{document}


\setmainfont[
    SmallCapsFont=*,
    ItalicFont=* Italic,
    SlantedFont=Helvetica Neue Light,
    BoldFont=Helvetica Neue Medium
]{Helvetica Neue Light}
\setsansfont[
    SmallCapsFont=*,
    ItalicFont=* Italic,
    SlantedFont=Helvetica Neue Light,
    BoldFont=Helvetica Neue Medium
]{Helvetica Neue Light}

%-----       resume       ---------------------------------------------------------
\makecvtitle

\section{Bildung}
\cventry{Seit Jul 2017}{Wissenschaftlicher Mitarbeiter/Ph.D.-Student}{\newline{}\href{http://hpi.de/}{Hasso-Plattner-Institut}}{Potsdam, Deutschland}{}{\textit{Forschungsschwerpunkt:} Integration von nahtlosem Lernen und mobilen Lerntechniken im Rahmen von Massive Open Online Courses}
\cventry{Okt 2013 \textendash\space Jun 2017}{Master of Science}{\href{http://hpi.de/}{Hasso-Plattner-Institut}}{Potsdam, Deutschland}{}{\textit{Studiengang:} IT-Systems Engineering;\newline{}
\textit{Masterarbeit:} Von MOOCs zu Micro Learning: Mobiles videobasiertes Lernen optimieren}
\cventry{Okt 2009 \textendash\space Jul 2012}{Bachelor of Science}{\href{http://hpi.de/}{Hasso-Plattner-Institut}}{Potsdam, Deutschland}{}{\textit{Studiengang:} IT-Systems Engineering;\newline{}
\textit{Bachelorarbeit:} Implementierung einer \href{https://apps.apple.com/de/app/seer-1000/id674468116}{iOS-Applikation}  zur Konfiguration von Langzeit-EKG- Rekordern über Bluetooth}
\cventry{Aug 2009}{Abitur}{Bohnstedt-Gymnasium}{Luckau, Deutschland}{}{}
%{Advanced courses: Mathematics and English}

\section{Erfahrung}
\cvexperience{Seit Jul 2017}{Software-Ingenieur}{\href{https://open.hpi.de}{openHPI}, \href{http://hpi.de/}{Hasso-Plattner-Institut}}{Potsdam, Deutschland}{
\begin{itemize}
	\item Hauptverantwortung für Konzeption, Entwicklung und Wartung einer nativen mobilen iOS-App für die HPI-Lernplattform \href{https://open.hpi.de}{\textit{openHPI}}
	\item Weiterentwicklung von \href{https://open.hpi.de}{\textit{openHPI}} mit dem Fokus auf öffentliche APIs und der Verbesserung von User Experience und Barrierefreiheit
	\item Leitung der Arbeitsgruppe für User Experience im \href{https://open.hpi.de}{\textit{openHPI}} -Entwicklerteam: Durchführung regelmäßiger Meetings sowie Planung und Umsetzung von Verbesserungsmaßnahmen
\end{itemize}
}
%\cvexperience{Mar 2016 \textendash\space Oct 2017}{Freelancer}{Potsdam, Deutschland}{}{Full Stack Software Developer \& Consultant.}
\cvexperience{Okt 2015 \textendash\space Jun 2017}{Lehrassistent}{\href{http://hpi.de/}{Hasso-Plattner-Institut}}{Potsdam, Deutschland}{Betreuung von Studententeams für den Kurs \href{https://hpi.de/uebernickel/teaching/global-team-based-innovation-gti-design-thinking.html}{Global Team-based Product Innovation (ME310)}}
\cvexperience{Mär 2014 \textendash\space Jun 2017}{Studentischer Mitarbeiter}{\href{http://hpi.de/}{Hasso-Plattner-Institut}}{Potsdam, Deutschland}{\textit{Lehrstuhl:} Enterprise Platform and Integration Concepts;\newline{}
Fußball-Analyse mit einer interaktiven Taktiktafel unter Verwendung von Geodaten}
\cvexperience{Apr 2013 \textendash\space Sep 2013}{Praktikum Software-Entwickler}{\href{http://www.sap.com/research}{SAP Research}}{Belfast, Vereinigtes Königreich}{\textit{Abteilung:} Next Business and Technologies;\newline{}
Entwicklung einer Geschäftsprozesssimulation auf der Grundlage von BPMN-Modellen}
\cvexperience{Aug 2012 \textendash\space Feb 2013}{Praktikum Mobile-Entwickler}{\href{http://www.getemed.net/en/}{Getemed AG}}{Teltow, Deutschland}{
Entwicklung einer \href{https://apps.apple.com/de/app/seer-1000/id674468116}{iOS-Applikation} zur Konfiguration von Langzeit-EKG-Rekordern}

\section{Technische Kompetenzen}
\cvitem{Spezialisierungen}{Software-Entwicklung, Mobile-Entwicklung, Datenanalyse, UI/UX-Design}
\cvitem{Qualifikationen}{Projektleitung, \href{https://www.scrum.org/user/1274563}{Scrum (PSM I)}, Design Thinking}
\cvitem{Interessen}{Adobe Illustrator, Adobe Photoshop}

\pagebreak

\subsection{Programmierkenntnisse}
\cvitem{Advanced}{Python, Swift, iOS-Entwicklung}
\cvitem{Intermediate}{Ruby, SQL, Backend-Entwicklung, Web-Entwicklung}
%\cvitem{Basic}{}

%\pagebreak

\section{Publikationen (Auszug)}
\cvpublication{Sep 2022}{\href{http://dx.doi.org/10.1109/LWMOOCS53067.2022.9927873}{Stay at Home and Learn: Did the COVID-19 Pandemic Influence Mobile\newline{}Learning in MOOCs?}}{M. Bothe, Ch. Meinel}{Proceedings of the  2022 IEEE Learning With MOOCS (LWMOOCS)}
\cvpublication{Nov 2021}{\href{http://dx.doi.org/10.1007/978-3-030-90677-1_5}{Video Consumption with Mobile Applications in a Global Enterprise MOOC Context}}{M. Bothe, F. Schwerer, Ch. Meinel}{Innovations in Learning and Technology for the Workplace and Higher Education (TLIC 2021)}
\cvpublication{Jun 2021}{\href{https://doi.org/110.1145/3430895.3460151}{The Impact of Mobile Learning on Students' Self-Test Behavior in MOOCs}}{M. Bothe, Ch. Meinel}{Eighth ACM Conference on Learning @ Scale 2021 (L@S '21)}
\cvpublication{Sep 2020}{\href{https://doi.org/10.1109/LWMOOCS50143.2020.9234368}{When Do Learners Rewatch Videos in MOOCs?}}{M. Bothe, Ch. Meinel}{Proceedings of the  2020 IEEE Learning With MOOCS (LWMOOCS)}
\cvpublication{Sep 2020}{\href{https://doi.org/10.1109/LWMOOCS50143.2020.9234338}{On the Potential of Automated Downloads for MOOC Content on Mobile\newline{}Devices}}{M. Bothe, Ch. Meinel}{Proceedings of the  2020 IEEE Learning With MOOCS (LWMOOCS)}
\cvpublication{Apr 2020}{\href{https://doi.org/10.1109/EDUCON45650.2020.9125145}{On the Acceptance and Effects of Recapping Self-Test Questions in MOOCs}}{M. Bothe, J. Renz, Ch. Meinel}{Proceedings of the 2020 IEEE Global Engineering Education Conference (EDUCON)}
\cvpublication{May 2019}{\href{https://doi.org/10.1007/978-3-030-19875-6_3}{Applied Mobile-Assisted Seamless Learning Techniques in MOOCs}}{M. Bothe, Ch. Meinel}{Digital Education: At the MOOC Crossroads Where the Interests of Academia and Business Converge}
\cvpublication{Apr 2019}{\href{https://doi.org/10.1109/EDUCON.2019.8725043}{From MOOCs to Micro Learning Activities}}{M. Bothe, J. Renz, T. Rohloff, Ch. Meinel}{Proceedings of the 2019 IEEE Global Engineering Education Conference (EDUCON)}
%\cvpublication{Mar 2019}{\href{https://www.researchgate.net/publication/331864553_Visualizing_Content_Exploration_Traces_of_MOOC_Students}{Visualizing Content Exploration Traces of MOOC Students}}{T. Rohloff, M. Bothe, Ch. Meinel}{Proceedings of the 9th International Conference on Learning Analytics \& Knowledge (LAK19)}
%\cvpublication{Sep 2018}{\href{https://doi.org/10.1109/LWMOOCS.2018.8534685}{Towards a Better Understanding of Mobile Learning in MOOCs}}{T. Rohloff, M. Bothe, J. Renz, Ch. Meinel}{Proceedings of the 5th Learning with MOOCs Conference (LWMOOCs)}
%\cvpublication{Nov 2017}{\href{https://doi.org/10.1145/3136907.3136931}{Supporting Multi-Device E-Learning Patterns with Second Screen Mobile Applications}}{T. Rohloff, J. Renz, M. Bothe, Ch. Meinel}{Proceedings of the 16th World Conference on Mobile and Contextual Learning (mLearn)}
%\cvpublication{Apr 2016}{\href{https://doi.org/10.5220/0005877600270035}{Recognizing Compound Events in Spatio-Temporal Football Data}}{K. Richly, M. Bothe, T. Rohloff, C. Schwarz}{Proceedings of International Conference on Internet of Things and Big Data (IoTBD)}

\section{Projekte (Auszug)}
%\cvproject{May 2016 - Jan 2018}{\href{https://crowdfunding.correctiv.org/strassengezwitscher/blog/}{Crowdgezwitscher}}{Development of an information platform for \href{https://twitter.com/streetcoverage}{Straßengezwitscher e.V.} reporting neutrally about xenophobic demonstrations.}
\cvproject{Jun 2017 - Aug 2024}{Mobile Anwendungen für die HPI-Lernplattform openHPI}{
Entwicklung von nativen mobilen Anwendungen für \href{https://github.com/openHPI/xikolo-ios}{iOS} und \href{https://github.com/openHPI/xikolo-android}{Android} für die HPI-Lernplattform openHPI, die ein übersichtliches Design, optimierte Interaktionen und ergänzende, auf eine mobile Nutzung zugeschnittene Funktionen bieten.}
\cvproject{Mär 2014 - Jun 2017}{Interaktive Taktiktafel für die Fußballanalyse}{Die Entwicklung einer interaktiven Taktiktafel ist eine wertvolle Hilfe bei der Spielbeurteilung und Trainingsvorbereitung, mit der Trainer und Profisportler analysieren können, wie ein Spiel aufgebaut ist und wie ein Teammitglied in einem Spiel agiert.}
\cvproject{Okt 2014 \textendash\space Jul 2015}{\href{https://hpi.de/uebernickel/teaching/global-team-based-innovation-gti-design-thinking.html}{Global Team-based Product Innovation \& Engineering (ME310)}}{Arbeit an einer Design-Herausforderung in Zusammenarbeit mit der Stanford University: Entwicklung eines Prototyps für die wichtigsten Merkmale eines Audi-Fahrzeugs im Jahr 2025 unter Anwendung von Design-Thinking-Methoden.}
\cvprojectwithsubtitle{Mär 2012 \textendash\space Feb 2013}{\href{https://apps.apple.com/de/app/seer-1000/id674468116}{SEER 1000 iOS-Applikation}}{Bachelor-Projekt \& Praktikum}{Entwicklung einer\href{https://apps.apple.com/de/app/seer-1000/id674468116}{iOS-Applikation} zur Konfiguration des Langzeit-EKG-Rekorders SEER 1000 über Bluetooth.}

\section{Weiteres}
\cvitem{Sprachen}{Deutsch (Muttersprache), Englisch (fließend)}
\cvitem{Sport}{Functional Fitness, Radfahren, Rudern}
\cvitem{Interessen}{Kochen, Familie, Reisen}

\end{document}
